\chapter{Methodology} \label{chap:method}


In this section, we explain the methoddology used in this thesis. 

Add a flow chart that summarizes the different steps in the pipeline, like the one in the interview presentation.

Add a table showing the cosmological parameters.

Consider adding a table which describes the fiducial cosmology.

\section{Modelling the electron power spectrum} \label{meth_sec:pee}

explain things in a schematic representation first. 

\subsection{The baryonification correction model}
explain how we ignore the redshift evolution of some parameters.
\subsection{The halo model}
Check that you have explained what dark matter halos are in the beginning. Refer to Either the review by Marika or the paper on the Web halo model.
Mention that we use the Tinker halo mass function and the NFW profile.

\subsubsection{The Core Cosmology Library (CCL)}



\section{Simulating the dispersion measure signal} \label{meth_sec:sim_dm}
Add a table with the fiducial values and their prior. explain how we ignore the redshift evolution of some parameters. Motivate the choice for the number of shells and the redshift range.

\subsection{The Generator for Large Scale Structure (GLASS)}
mention the reason behind the choice of 30 shells.

\subsection{The electron angular power spectrum}
\subsection{The cosmological covariance of the dispersion measure signal}


\section{Simulating the systematic uncertainties}

\subsection{Milkyway contribution}
\subsection{Host contribution}

\subsection{Redshift uncertainty}
I simulate the imapct of redshift uncertainty by assuming that a redshift measuremnet $z_i$ is sampled from a Gaussian distribution with mean $z_i$ and standard deviation $\sigma_i = \eta z_i$, where $\eta_i$ is a number between $0$ and $1$. To propagate the error on the cosmological covariance, consider that $(z_i, z_j)  \sim \mathcal{N}((\hat{z}_i, \hat{z}_j) , \Sigma)$, with $\mathcal{N}$ is the normal distribution and
\begin{equation}
    \Sigma = \begin{bmatrix}
    \sigma_i^2 & 0 \\
    0 & \sigma_j^2 
    \end{bmatrix}
\end{equation}
To propagate the error, I calcualte the expectation value $\langle \hat{C}_{ij} \rangle$, given as
\begin{equation}
    \langle \hat{C}_{ij} \rangle = \int \mathcal{N}((\hat{z}_i, \hat{z}_j) , \Sigma) \, \, C_{ij} \, \,dz_i \, \, dz_j
\end{equation}
With Taylor's expansion around $(\hat{z}_i, \hat{z}_j)$ and setting $\Delta_{ij} = (z_i, z_j) - (\hat{z}_i, \hat{z}_j)$, the covariance matrix can be written
\begin{equation}
    \langle \hat{C}_{ij} \rangle = \int \mathcal{N} \, \, \left(\hat{C}_{ij} + \nabla C_{ij}|_{(\hat{z_i}, \hat{z_j})} \Delta_{ij} + \frac{1}{2} \Delta_{ij}^{T} \nabla^2 C_{ij}|_{(\hat{z_i}, \hat{z_j})}\Delta_{ij} + .....\right)\, \,dz_i \, \, dz_j.
\end{equation}
Since the Gaussian distribution has vanaishing moments for moment larger than third order, the expression is simplified into
\begin{equation}
    \langle \hat{C}_{ij} \rangle = \hat{C}_{ij} + \frac{1}{2} \, \, \sigma_{i}^2 \, \, \partial^2_{ii} \, C_{ij}|_{(\hat{z_i}, \hat{z_j})} \, \, + \frac{1}{2} \, \, \sigma_{j}^2 \, \, \partial^2_{jj} \, C_{ij}|_{(\hat{z_i}, \hat{z_j})}.
\end{equation}
Therefore, it is possible to account for the error in the covariance due to redshift uncertainty using the second derivative of the covariance with respect to redshift $z$.

I show the impact of redshift uncertainty on the scatter in the DM-$z$ relation in chapter \ref{chap:results}.


\subsection{Selection effects}
\subsection{Unlocalised hosts}




\section{Simulation-based inference} \label{meth_sec:sbi}
\subsection{Approximate Baysian inference}
\subsection{Neural Posterior Estimator}
\subsection{The Sequentual Neural Posterio Estimator}
\subsubsection{The truncated sequential neural posterior estimator}
\subsection{Flow matching posterior estimator}

\subsection{Data Compression}





I simulate the electron density shells and place FRBs in the same redshift and angular position as in the mock catalog. This will reduce the computational cost and the need to calculate the covariance matrix for each realisation within the SBI pipeline. (Try to do this excercise without choosing the sdame redshifts for a lower number of sources such that it lowers down the computational costs)
\subsubsection{Numerical derivative stability}
\subsubsection{The Sensitivity of compressed statistics to model parameters}

Add a plot about derivative stability. Add a plot for the sensitivity.


